\section{First Collapse: A Skill Loses Its Social Position}

A person can remain employed while the meaning of the job changes beneath them. The organization still wants a name on the workflow. It needs less of the craft through which that person learned to justify the name.

The first collapse is a loosening of the connection among capability, output, recognition, and position.

\subsection{The Interface Changes Before the Job Disappears}

A professional report once signaled more than a collection of sentences. It suggested access to information, trained judgment, command of a field, and time spent on the work. When a similar-looking report becomes cheap to produce, the signal changes.

The output may improve. The person's actual competence may remain. What weakens is the output's ability to distinguish the person who produced it.

A simple signal-detection expression makes the idea visible. For two Gaussian signal distributions with equal standard deviation $\sigma$,
\begin{equation}
 d'=\frac{\mu_1-\mu_0}{\sigma}.
\end{equation}
If the relevant distributions overlap more, discrimination becomes harder. This is an analogy for changing professional signals, not an estimate of AI's labor-market effect. A polished document now needs other evidence if it is meant to certify individual judgment.

Education has encountered related positional dynamics. Collins's account of credential competition examines how qualifications can become thresholds when their supply expands faster than the positions they distinguish \cite{collins1979credential}. The knowledge can remain useful while its exchange value changes.

Machine assistance can also complement a skill and raise its return. A capable accountant may serve more clients. An engineer may investigate designs previously beyond the budget. A novice may learn faster. The question is which complementary positions organizations actually create and who can reach them.

The familiar ATM case shows why task substitution and employment are different variables. Bessen's historical analysis describes how automation reduced teller work per branch while changes in branch networks and customer-facing work supported other demand \cite{bessen2015learning}. It supplies a mechanism, rather than a guarantee that each new automation wave will reproduce the same result.

Current evidence also needs this separation. The Stanford Digital Economy Lab's August 2026 revision, using payroll data through June, reports no widespread economy-wide displacement while documenting a substantial divergence for younger workers in exposed occupations \cite{brynjolfsson2026canaries}. The authors treat the findings as descriptive indicators, with reduced hiring an important channel. Aggregate stability can coexist with a difficult entry path.

That entry path matters beyond the immediate job count. Routine junior work can be part of how people acquire judgment. If automation removes the practice while retaining a demand for experienced reviewers, organizations need another way to produce the experience.

\subsection{Having the Result and Having Done the Work}

The second separation is personal. Receiving a finished result differs from participating in the process that made it possible.

Arendt distinguished labor, work, and action: recurring life-maintenance, durable making, and activity among plural people in a public world \cite{arendt1958human}. The distinction explains why a job can provide several kinds of attachment at once. It supplies income, an object produced, and a recognizable place among others.

Automating production can loosen those attachments together. Telling someone to “learn the tool” may solve the immediate delivery problem while leaving the question of recognition unanswered.

Kaczynski's account of a “power process” describes satisfaction through forming a goal, exerting effort, overcoming resistance, and achieving it \cite{kaczynski1995industrial}. The sequence is useful as a description of experienced agency. His larger claims about technological inevitability require their own arguments. The local mechanism does not establish a necessary political or historical endpoint.

In the book's coordinates, delegated work can make the user's immediate $R_p$, $T_p$, and $G_p$ small:
\begin{equation}
 \eta_p=\eta_0
 \exp(\alpha R_p+\beta T_p+\gamma G_p)
 \longrightarrow\eta_0
 \quad\text{as }(R_p,T_p,G_p)\to(0,0,0).
\end{equation}
The limit is the baseline $\eta_0$. More importantly, it concerns the user's experienced portion of a task. The provider and receiving world still have bills. This illustrative index does not quantify fulfillment.

The distinction is familiar outside AI. Paying for a meal and learning to cook both meet a need, but they produce different capabilities. Delegation can free time for something more meaningful. It can also remove an activity through which someone maintained confidence and competence. The resulting life depends on what fills the space.

Learned-helplessness research concerns controllability under specific experimental conditions \cite{seligman1967failure,maier2016learned}. It offers a useful question about whether action changes outcomes. It cannot be transferred wholesale into a diagnosis of people who use automated tools.

The stronger social question is whether people retain places where their judgments have intelligible effects. An environment that supplies endless results while rarely needing a person's sustained contribution can weaken that connection, even when its products are excellent.

\subsection{Demand Can Exist Without a Funded Position}

Care, education, repair, public service, and community life contain abundant unmet needs. Those needs do not automatically create paid roles. Budgets, procurement, professional rules, bargaining power, and institutions connect a useful activity to a livelihood.

Baumol's analysis explains how differing productivity growth can raise the relative cost of labor-intensive services \cite{baumol1967macro}. A society can become more productive overall while struggling to finance activities that still need time and presence.

This is why “people will move to more human work” is incomplete. The work needs an organization able and willing to pay, train, authorize, and retain the people doing it.

Arendt's analysis of superfluousness addresses a much heavier historical and political setting \cite{arendt1951origins}. The structural echo is the loss of an effective place in a shared world. It should sharpen attention to institutions, rather than turn ordinary workplace change into an automatic analogy with totalitarianism.

A person who loses an old occupational signal has not lost human worth. The economic interface has changed. Treating this solely as a motivation problem asks an individual to repair a relation that the surrounding system no longer offers on the old terms.

The first collapse can therefore occur without a layoff notice. People may still produce more, earn a salary, and use increasingly capable tools while becoming less certain where their judgment enters. A new position requires a real connection among contribution, learning, recognition, and authority.

When those connections remain weak, immediate feedback elsewhere becomes more attractive. That opens the second route of the chapter.
